\scanpage{12}
% archive foliation: 11
 i.e la donnée d'une situation $X'_i\mathbin{-}X_i$ sur $k_i$, \deleted{dont} et un iso. $(\mathcal X'\mathbin{-}\mathcal X)\simeq(X'_i\mathbin{-}X_i)\otimes_{k_i}\unsure{\bar k}$; de ceci se déduit une extension de schémas en groupes finis étales sur $k_i$
\[
(27)\quad\diagram{\begin{tikzcd}[ampersand replacement=\&]
1\ar[r,no head]\&G_i\ar[r,no head]\&\mathbb G'_i\ar[r]\&\mathbb G_i\ar[r,no head]\&1
\end{tikzcd}}
\]
et on exige que $\mathbb G'_i$ soit un groupe constant [i.e. que les \unsure{autom.} ``géom'' de $\mathcal X'\mathbin{-}\mathcal X$ (sur $\bar k$) \deleted{soient} ``définis sur $k_i$'' $\ldots$]

Il faut bien noter que pour $k_i$ fixé, il peut y avoir plusieurs données non isomorphes de cette espèce sur $k_i$, i.e. plusieurs scindages partiels $\Gamma_i\mathbin{-}\mathcal E'$ qui centralisent $\mathbb G'$. \deleted{\unsure{La}} La propriété d'unicité \deleted{consiste} évidente consiste pour $\unsure{\mathrm{tt}|\forall}$ \unsure{s'égaliser|s'égaler}, quitte à passer à une \unsure{surextension} finie assez grande $\ldots$

De façon précise, introduisons
\[
(28)\qquad\unsure{\mathcal Z|Z}=\operatorname{Centr}_{\mathcal G'}(\mathbb G')
\]
